\chapter{Liver conditions and diseases}

\section*{Definition and Overview}
Liver conditions and diseases encompass a broad spectrum of pathological conditions affecting the hepatobiliary system, which performs critical metabolic, synthetic, detoxicating, and immunological functions. The liver is susceptible to infectious, toxic, metabolic, autoimmune, vascular, genetic, and neoplastic insults. Common conditions include metabolic dysfunction-associated steatotic liver disease (formerly non-alcoholic fatty liver disease), alcoholic liver disease, viral hepatitis (A, B, C, D, E), autoimmune hepatitis, hereditary disorders (such as haemochromatosis and Wilson's disease), biliary diseases (such as primary biliary cholangitis and primary sclerosing cholangitis), and primary liver cancer (hepatocellular carcinoma). Left untreated, persistent liver inflammation (hepatitis) leads to progressive parenchymal destruction, fibrotic scarring (cirrhosis), and ultimate hepatic insufficiency or portal hypertension, necessitating liver transplantation.

\section*{Epidemiology}
Liver disease is a major cause of morbidity and mortality worldwide, with a rising prevalence largely driven by epidemics of obesity, diabetes, and alcohol misuse. In Australia, chronic liver disease affects millions of individuals, and hepatocellular carcinoma is one of the fastest-growing causes of cancer-related mortality. Metabolic dysfunction-associated steatotic liver disease is the most common liver condition globally, affecting up to 25--30\% of the general population in Western countries, mirroring rates of obesity and type 2 diabetes. Alcoholic liver disease remains a leading cause of cirrhosis and liver transplantation. Viral hepatitis, particularly chronic hepatitis B and C, affects hundreds of thousands of Australians; while hepatitis C prevalence has declined due to direct-acting antiviral therapies, hepatitis B remains a significant health burden, particularly in migrant and indigenous populations.

\section*{Aetiology and Pathophysiology}
The aetiologies of liver disease are diverse, but they share common pathways of hepatocyte injury, inflammatory response, and fibrogenesis. Persistent injury triggers the activation of hepatic stellate cells, which transition into myofibroblasts and deposit excessive extracellular matrix (collagen), leading to progressive fibrosis and eventual cirrhosis.
\begin{itemize}
    \item \textbf{Metabolic dysfunction:} Insulin resistance leads to excess accumulation of triglycerides within hepatocytes (steatosis), triggering oxidative stress, lipotoxicity, and inflammatory progression to steatohepatitis.
    \item \textbf{Toxins and alcohol:} Chronic ethanol consumption produces acetaldehyde, a highly reactive toxic metabolite that causes direct mitochondrial damage, lipid peroxidation, and cell necrosis.
    \item \textbf{Viral infection:} Hepatotropic viruses cause liver injury. While Hepatitis A and E cause acute infection, Hepatitis B and C can establish chronic infections, driving immune-mediated destruction of hepatocytes over decades.
    \item \textbf{Autoimmune mechanisms:} Loss of self-tolerance leads to immune-mediated attack on hepatocytes (autoimmune hepatitis) or the biliary epithelial cells (primary biliary cholangitis and primary sclerosing cholangitis).
    \item \textbf{Genetic mutations:} Inborn errors of metabolism, such as \textit{HFE} gene mutations causing hereditary haemochromatosis (excess iron deposition) or \textit{ATP7B} mutations causing Wilson's disease (excess copper accumulation).
\end{itemize}

\section*{Clinical Features}
Early-stage liver disease is notoriously asymptomatic, with abnormalities often detected incidentally on routine blood tests. Symptomatic liver disease typically reflects advanced fibrosis, cirrhosis, or acute liver failure.
\begin{itemize}
    \item Nonspecific symptoms such as chronic fatigue, lethargy, mild right upper quadrant abdominal discomfort, and loss of appetite.
    \item Jaundice (yellowing of the skin and sclera) and dark-coloured urine, caused by impaired bilirubin conjugation or excretion.
    \item Pruritus (intense skin itching) resulting from the accumulation of bile acids or other pruritogens in the skin.
    \item Ascites (fluid accumulation in the peritoneal cavity) and peripheral oedema (swelling of the ankles and legs) due to portal hypertension and hypoalbuminaemia.
    \item Easy bruising, petechiae, and prolonged bleeding from minor cuts, due to impaired synthesis of clotting factors.
    \item Hepatic encephalopathy, characterised by confusion, cognitive decline, sleep-wake reversal, and asterixis (flapping tremor), caused by accumulation of ammonia and other neurotoxins.
    \item Cutaneous signs including spider naevi (dilated blood vessels on the chest and face), palmar erythema (reddening of the palms), and gynaecomastia or testicular atrophy in men due to impaired oestrogen clearance.
    \item Haematemesis or melaena from ruptured oesophageal varices, which are dilated veins resulting from portal hypertension.
\end{itemize}

\section*{Differential Diagnosis}
The clinical features of liver disease can mimic or be caused by other systemic conditions:
\begin{itemize}
    \item \textbf{Congestive heart failure:} Can cause passive hepatic congestion ("cardiac cirrhosis"), resulting in hepatomegaly, ascites, and abnormal liver enzymes.
    \item \textbf{Haemolytic anaemia:} Causes pre-hepatic jaundice and elevated unconjugated bilirubin, but is distinguished by a low haemoglobin, reticulocytosis, and normal liver function tests.
    \item \textbf{Biliary tract obstruction:} Choledocholithiasis (common bile duct stones) or pancreatic cancer can cause post-hepatic jaundice, pruritus, and elevated alkaline phosphatase, mimicking intrahepatic cholestasis.
    \item \textbf{Drug-induced liver injury:} Acetaminophen (paracetamol) overdose or idiosyncratic reactions to medications (such as amoxicillin-clavulanate) or herbal supplements can mimic acute viral or autoimmune hepatitis.
    \item \textbf{Sepsis and systemic infection:} Can cause secondary cholestasis and liver dysfunction, mimicking primary liver disease.
\end{itemize}

\section*{When to Seek Medical Care}
Individuals with risk factors, such as heavy alcohol use, obesity, or a family history of liver disease, should discuss liver health with their general practitioner. Anyone showing signs of jaundice or abdominal swelling should seek prompt medical review.

\textbf{Call 000 (or local emergency services) for:}
\begin{itemize}
    \item Vomiting blood (haematemesis) or passing black, tarry stools (melaena).
    \item Sudden onset of severe confusion, extreme drowsiness, or altered consciousness.
    \item Severe, acute abdominal pain, particularly if accompanied by high fever (possible spontaneous bacterial peritonitis).
    \item Difficulty breathing or chest tightness associated with severe abdominal swelling.
\end{itemize}

Other non-emergency reasons to seek medical care include:
\begin{itemize}
    \item New yellowing of the eyes or skin (jaundice).
    \item Unexplained swelling in the abdomen (ascites) or ankles.
    \item Persistent, severe itching without a skin rash.
    \item Unexplained fatigue, weight loss, or abdominal fullness.
\end{itemize}

\section*{Diagnosis and Investigations}
Diagnostic evaluation of liver disease combines blood tests, non-invasive fibrosis assessment, imaging, and occasionally a liver biopsy.
\begin{itemize}
    \item \textbf{Liver function tests:} A panel of blood tests measuring serum transaminases (ALT, AST, indicating hepatocyte injury), alkaline phosphatase and gamma-glutamyl transferase (indicating cholestasis), bilirubin, and albumin.
    \item \textbf{Coagulation profile:} Measures the international normalised ratio (INR) to assess the liver's synthetic capacity for clotting factors.
    \item \textbf{Viral serology:} Tests for Hepatitis A, B, C, and E antigens, antibodies, or viral PCR to detect active or past infection.
    \item \textbf{Autoimmune markers:} Anti-nuclear antibody, anti-smooth muscle antibody, and anti-mitochondrial antibody to screen for autoimmune hepatitis or primary biliary cholangitis.
    \item \textbf{Metabolic and genetic screening:} Serum ferritin and transferrin saturation (for haemochromatosis), caeruloplasmin (for Wilson's disease), and alpha-1 antitrypsin levels.
    \item \textbf{Abdominal ultrasound:} The primary imaging modality, which can detect fatty infiltration, gallstones, biliary dilatation, cirrhosis, splenomegaly, ascites, and focal liver masses.
    \item \textbf{Transient elastography:} A non-invasive ultrasound-based test (FibroScan) that measures liver stiffness to estimate the degree of hepatic fibrosis.
    \item \textbf{Liver biopsy:} The definitive diagnostic gold standard, used to grade inflammation, stage fibrosis, and identify specific aetiological features when non-invasive tests are inconclusive.
\end{itemize}

\section*{Management}
The management of liver disease is tailored to the specific underlying cause, with the primary goals of stopping disease progression, managing complications, and preventing liver failure.
\begin{itemize}
    \item \textbf{Lifestyle modifications:} For metabolic liver disease, weight loss of 7--10\% via a healthy diet (e.g., Mediterranean diet) and regular cardiovascular exercise is the mainstay of treatment.
    \item \textbf{Alcohol cessation:} Absolute abstinence is crucial for all stages of alcoholic liver disease and helps reverse early steatosis and slow the progression of cirrhosis.
    \item \textbf{Antiviral therapy:} Direct-acting antiviral therapy for Hepatitis C (highly effective, curable in 8--12 weeks) and nucleoside/nucleotide analogues (e.g., tenofovir, entecavir) to suppress Hepatitis B replication.
    \item \textbf{Immunosuppressive therapy:} Corticosteroids (e.g., prednisolone) and azathioprine to suppress the immune system in autoimmune hepatitis.
    \item \textbf{Chelation and venesection:} Therapeutic venesection (phlebotomy) to remove excess iron in haemochromatosis, and copper chelating agents (e.g., penicillamine) for Wilson's disease.
    \item \textbf{Management of cirrhosis complications:} Low-sodium diet and diuretics (spironolactone, furosemide) for ascites; beta-blockers (propranolol) or endoscopic band ligation for oesophageal varices; lactulose and rifaximin for hepatic encephalopathy.
    \item \textbf{Liver transplantation:} The final curative option for patients with end-stage liver disease (decompensated cirrhosis) or acute liver failure who meet specific transplant criteria.
    \item \textbf{Hepatocellular carcinoma screening:} Six-monthly abdominal ultrasound and serum alpha-fetoprotein measurements in patients with cirrhosis or chronic hepatitis B.
\end{itemize}

\section*{Complications}
\begin{itemize}
    \item \textbf{Decompensated cirrhosis:} Characterised by hepatic encephalopathy, variceal haemorrhage, ascites, and jaundice.
    \item \textbf{Spontaneous bacterial peritonitis:} A life-threatening bacterial infection of ascites fluid, requiring urgent antibiotics.
    \item \textbf{Hepatorenal syndrome:} Functional renal failure occurring in patients with advanced liver disease, caused by renal vasoconstriction.
    \item \textbf{Hepatopulmonary syndrome:} Pulmonary vascular dilatation leading to hypoxaemia (low oxygen levels) in cirrhotic patients.
    \item \textbf{Hepatocellular carcinoma:} Primary liver cancer that almost exclusively arises in the setting of chronic liver inflammation and cirrhosis.
    \item \textbf{Acute liver failure:} Rapid loss of hepatic function (within days or weeks) in a person without pre-existing liver disease, most commonly caused by paracetamol overdose or acute viral hepatitis.
\end{itemize}

\section*{Prognosis and Outcomes}
The prognosis of liver disease depends heavily on the specific aetiology, the stage of fibrosis at the time of diagnosis, and the patient's adherence to management. Early-stage conditions like simple steatosis or early fibrosis are highly reversible with lifestyle changes or eradication of the underlying cause (e.g., viral clearance). Once cirrhosis develops, the disease is generally irreversible, but it can be managed in a compensated state for many years. However, once decompensation occurs (development of ascites, variceal bleeding, or encephalopathy), the prognosis drops significantly, with a 5-year survival rate of less than 50\% without a liver transplant. Liver transplantation offers excellent outcomes, with 5-year survival rates exceeding 80--85\%.

\section*{Prevention and Self-Care}
\begin{itemize}
    \item Limit alcohol consumption to recommended safe guidelines, or abstain completely if you have pre-existing liver disease.
    \item Maintain a healthy body weight through a balanced, low-sugar diet and regular exercise to prevent metabolic-associated fatty liver disease.
    \item Get vaccinated against Hepatitis A and Hepatitis B, especially if you are in a high-risk group or planning to travel.
    \item Practice safe sex (using barrier methods) and avoid sharing needles or personal items like razors to prevent the transmission of Hepatitis B and C.
    \item Use medications, including over-the-counter paracetamol, strictly according to the package directions, and avoid combining them with alcohol.
    \item Avoid unregulated herbal supplements and traditional medicines, which are common causes of drug-induced liver injury.
    \item Attend all recommended medical appointments and routine blood tests if you have been diagnosed with chronic liver disease.
\end{itemize}

\section*{Sources and Further Reading}
The following organisations provide authoritative, evidence-based information on liver diseases, support, and resources:
\begin{itemize}
    \item Liver Foundation. \textit{Support, education, and resources for liver health in Australia.} https://liver.org.au/
    \item Gastroenterological Society of Australia. \textit{Clinical guidelines and patient information on liver diseases.} https://www.gesa.org.au/
    \item British Liver Trust. \textit{Authoritative patient guides on liver conditions and prevention.} https://britishlivertrust.org.uk/
    \item American Association for the Study of Liver Diseases. \textit{Professional reference and patient information.} https://www.aasld.org/
    \item World Health Organization. \textit{Global hepatitis report and prevention strategies.} https://www.who.int/
\end{itemize}
